\begin{frame}{Rule 1: Efficiency --- When to Reduce Delay vs. Compensate}
    \begin{center}
        \textbf{Choose the policy that achieves uptake at the lowest \textbf{cost}.}
    \end{center}

    \vspace{1.5em}
    \begin{center}
        \Large
        Reduce \textbf{Delay} if: \textbf{Cost} of speed-up $<$ \textbf{MWTA} \\
        \textbf{Compensate} if: \textbf{Cost} of speed-up $>$ \textbf{MWTA}
    \end{center}
    \vspace{1.5em}

    \begin{columns}[T]
        \begin{column}{0.48\textwidth}
            \textbf{Intuition}
            \begin{itemize}
                \item \textbf{MWTA} measures how costly delay is to individuals
                \item Speed-up reduces delay directly
                \item Compensation offsets delay without changing timing
            \end{itemize}
        \end{column}
        
        \hfill
        
        \begin{column}{0.48\textwidth}
            \textbf{Implication}
            \begin{itemize}
                \item If delay is very costly, reducing it is more \textbf{efficient}
                \item If speed-up is expensive, incentives may be preferable
                \item Policy should compare marginal \textbf{costs} on both sides
            \end{itemize}
        \end{column}
    \end{columns}
    
    \vfill
    \begin{center}
        \textit{Efficiency requires matching policy choice to the cost of delay.}
    \end{center}
\end{frame}
